\documentclass[10pt,a4paper]{amsart}
\usepackage[margin=19mm]{geometry}
\usepackage{amsmath,amssymb,mathtools}
\usepackage[scr=rsfs,cal=euler]{mathalfa}
\usepackage{xfrac}
\usepackage[unicode,hidelinks,bookmarks=false]{hyperref}
\newcommand{\ie}{i.e.\ }
\newcommand{\cf}{cf.\ }
\newcommand{\resp}{respectively\ }
\newcommand{\Subsection}{\subsection}
\providecommand{\nobreakdash}{\nobreak\hbox{-}}
\setlength{\emergencystretch}{2em}
\multlinegap0pt
\usepackage{float}
\usepackage{amssymb}
\usepackage[all]{xy}
\usepackage[normalem]{ulem}
\usepackage{tcolorbox}
\usepackage{xfrac}
\usepackage{caption}
\newtheorem{theorem}{Theorem}
\usepackage{cleveref}
\crefname{theorem}{Theorem}{Theorems}
\title{Graphs arising from the dual Steenrod algebra}
\author{Connor Elliott, Courtney Hauf, Kai Morton, Sarah Petersen, Leticia Schow}
\date{Source: arXiv:2508.17041v2 (CC BY 4.0)}
\begin{document}
\maketitle

\noindent\emph{Source and licence.} Graphs arising from the dual Steenrod algebra. Version arXiv:2508.17041v2 (CC BY 4.0), distributed by arXiv under the Creative Commons Attribution 4.0 International licence, http://creativecommons.org/licenses/by/4.0/, as recorded in the arXiv licence metadata for this exact version and shown on the abstract page https://arxiv.org/abs/2508.17041v2. Full licence notice: \texttt{licenses/arxiv-2508.17041-CC-BY-4.0.txt}.\par\smallskip

\noindent\textit{Editorial scope.} Counted unit: the article's theorem thm:AntipodeOdd (source0.tex lines 1414--1418). The article's complete original proof of it is delivered with it (source0.tex lines 1420--1443, one \texttt{proof} environment of 24 lines). No mathematics is written, repaired or re-derived: the statement and the proof are the article's own lines, in the article's own notation, and no retained line is substituted or re-flowed.  No further source-local context is needed: every cross-reference of every retained line resolves inside the two delivered spans.  Nothing was omitted from any retained span, so the package carries no omission record.  No retained line contains a citation, so the wrapper carries no bibliography and no bibliography entry.  No retained line contains a figure environment, an image-inclusion command or drawing code, so no image is delivered and the package declares no asset dependency.  Editorial formatting only, in addition: an AMS \texttt{amsart} preamble that restates the article's own declarations and macros used by the retained lines, namely the environments \texttt{theorem} on separate counters, together with the standard packages those lines need.  The retained environments print in this document as Theorem 1 in order of appearance, whereas the article prints the same items with the numbering of its own full document.  The complete native source of the article as distributed in the arXiv e-print for this exact version is bundled unchanged as \texttt{sources/183743/source0.tex}.
\begin{theorem} \label{thm:AntipodeOdd} 
    The antipode $c(\xi_i^{p^j}) \in A^\vee(n)$ is the signed sum of all directed paths from the vertex $2 p^j$ to the vertex $2 p^{i + j}$ in the complete graph on the vertices $\{1,2, 2p, 2p^2, \cdots, 2p^n\}$ considered as a directed graph. The positive terms in the sum correspond exactly to paths consisting of an even number of edges from the vertex $2$ to the vertex $2p^{i +j}$.

    Similarly, the antipode $c(\tau_i) \in A^\vee(n)$ is the signed sum of all directed paths from the vertex $1$ to the vertex $2 p^{i}$. The positive terms correspond exactly to paths consisting of an even number of edges from the vertex $1$ to the vertex $2 p^i$. \color{black}
\end{theorem} 

\begin{proof}
Recall that the antipode map $c: A^\vee (n) \to A^\vee (n)$ of the Hopf algebra structure sends
\begin{align*}
    c(\xi_n) & = - \xi_n - \sum_{k = 0}^{n - 1} \xi_{n - k}^{p^k} c (\xi_k),
\end{align*}
and
\begin{align*}
    c(\tau_n) & = - \tau_n - \sum_{k = 0}^{ n - 1} \xi_{n - k}^{p^k} c(\tau_k).
\end{align*}
We will prove the statement for $c(\xi_i^{p^j})$ by induction on $i$. The proof for $c(\tau_i)$ is similar. 

To begin the induction, consider $c(\xi_1^{p^j}) = c(\xi_1)^{p^j} = (-\xi_1)^{p^j} = - \xi_1^{p^j}$ since $p$ is an odd prime. In the  complete graph on the vertices $\{1,2, 2p , 2p^2, \cdots, 2p^n\}$, the monomial $\xi_1^{p^j}$ corresponds exactly to the a path from $2p^j$ to $2p^{j + 1}$.\color{black} This is a path of length one so we observe the desired statement is true when $i = 1$. 

Now suppose $c(\xi_\ell^{p^j})$ is the signed sum of all directed paths from $2{p^j}$ to $2 p^{\ell + j}$ for all $\ell \leq i$ where paths of an even length correspond to positive terms and paths of an odd length correspond to negative terms. We must show that 
\begin{equation} \label{eq:conj}
c(\xi_{i + 1}^{p^j}) = - \xi_{i + 1}^{p^j} - \sum_{k = 0}^{i} \xi_{i + 1 - k}^{p^{k + j}} c(\xi_k^{p^j})
\end{equation}
is the signed sum of all directed paths from $2p^j$ to $2p^{i + 1 + j}$ where directed paths of an even length correspond to positive terms and directed paths of an odd length correspond to negative terms.

Consider the term $- \xi_{i + 1 - k}^{p^{k + j}} c(\xi_k^{p^j})$ in \cref{eq:conj}. The $c(\xi_k^{p^j})$ factor in this term consists of a signed sum of all paths from $2 p^j$ to $2 p^{j + k}$ while the $\xi_{i + 1 - k}^{p^{k + j}}$ factor consists of the edge from $2 p^{j + k}$ to $2 p^{j + i +1}$. Thus $\xi_{i + 1 - k}^{p^{k + j}} c(\xi_k^{p^j})$ is a signed sum of all directed paths from $2 p^j$ to $2 p^{j + k}$ followed by the edge from $2 p^{j + k}$ to $2 p^{j + i +1}$ in the reduced graph $G_x^\textup{R}$. Observe that the length of each directed path in $-\xi_{i + 1 - k}^{p^{k + j}} c(\xi_k^{p^j})$ is one step longer than in $c(\xi_k^{p^j})$. Moreover, the sign in front of each term appearing in $-\xi_{i + 1 - k}^{p^{k + j}} c(\xi_k^{p^j})$ is the opposite of that appearing in $c(\xi_k^{p^j})$. Hence directed paths of even length still correspond to positive terms in the sum and directed paths of odd length still correspond to negative terms in the sum.
Thus summing over all $k$ where $0 \leq k \leq i$ gives the desired signed sum of all paths from $2 p^j$ to $2 p^{i + 1 + j}$ passing through at least one intermediate vertex. The term $-\xi_{i + 1}^{p^j}$ gives the remaining path from $2 p^j$ to $2 p^{i + 1 +j}$ passing through no intermediate vertices. Hence $c(\xi_{i + 1}^{p^j})$ is the signed sum of all paths from $2p^j$ to $2 p^{i + 1 + j}$ where directed paths of an even length correspond to positive terms and directed paths of an odd length correspond to negative terms. 


\end{proof}

\end{document}
